\chapter{Block models}\label{ch:blocks}
\begin{goals}
\textbf{In\index{block model} this chapter} you open regular, sub-blocked and rotated block models from CSV, Datamine and Parquet files, learn how micro finds the lattice and what it does when there is none, choose a format to keep a model in, and contour a grade shell.
\end{goals}

To micro a block model is a table with one row per block and three columns that hold the block's centroid. Everything else (the block size, the origin, the counts, whether the model is sub-blocked or rotated) micro works out from the centroids themselves. You do not type a block size. You do not need a definition file beside the CSV.

When micro finds the lattice, it draws every block as a true box, cuts it true on a section, picks it exactly, and offers the analyses that need a volume. When it finds none, the model still opens, as points. Points have no volume, and so no tonnage. This chapter is about the difference, and about how to get from one to the other.

\section{How micro finds the lattice}\label{sec:blocks-lattice}

Opening\index{lattice} a delimited model takes two passes over the file. The first reads the three coordinate columns and looks, on each axis, for the spacing that every centroid is a whole number of steps along: that is the block size on that axis, and the smallest centroid is the origin. The second pass streams the blocks to the screen, each placed by its integer position on that lattice.

For the first pass micro needs to know which columns are the coordinates. It recognizes the usual names, and the open dialog lets you correct the choice (\cref{sec:opening-dialog}). It also picks one numeric column to color by and, if there is one, a text column to use for categories; you can change both afterwards (\cref{ch:color}).

The result is on the status line and in file info (\key{i}): the block size on each axis, the counts, and the number of blocks present. A model need not fill its box. The demo's \file{deposit.csv} has no blocks above the pit surface, and missing blocks cost nothing.

\begin{note}
A row with a blank coordinate is not drawn, since it has nowhere to go, and is not counted as a record by any analysis. A blank is never read as zero (\cref{sec:columns-blanks}).
\end{note}

\subsection{When there is no lattice}

Some files are not block models, or not regular ones: a model exported with rounded coordinates, a cloud of estimation points, a model whose sub-blocks follow no simple rule. micro does not force a grid onto them. The layer opens as \term{points}: one dot per row, colored, filterable and pickable like any model, with statistics and export available. Grade–tonnage, swath, the cross-tab and grade shells are not offered, because a dot has no volume.

You can also choose points deliberately, for a quick look at a huge model or when you suspect the detected grid is wrong: \ui{Open as} \ui{points} in the open dialog, or \menu{Data \& file>Reinterpret as…>points} on the layer's menu (\cref{sec:opening-openas,sec:opening-reinterpret}).

\section{Looking at a model}\label{sec:blocks-display}

Blocks are drawn as boxes with their real faces, so a model reads as a solid and a thin section reads as a painted wall. Four display choices belong to the model, not to its colors:

\begin{itemize}
\item \textbf{Block edges} outline every block, and every cut block on a section. \sidekey{⇧B}{block edges on or off, for every model} \key{⇧B} toggles them for the whole view; \menu{Display>Block edges} on a layer's right-click menu overrides that for one model (\ui{auto}, \ui{on}, \ui{off}).
\item \textbf{Opacity} (\menu{Display>Opacity} on the layer's menu: 100, 75, 50 or 25\,\%) makes a model see-through, for a shell or drillholes inside it.
\item \textbf{Boxes or points} for all models at once is in the \ui{View} menu: \ui{Blocks as points} draws fixed splats and is snappier on a weak graphics card; \ui{Blocks as boxes} returns to true faces.
\item \textbf{Point size} (×1 to ×8) appears on the menu of a model drawn as points.
\end{itemize}

\Cref{ch:color} covers color, ramps, classes and the legend; \cref{ch:sections} covers sections.

\begin{tryit}
\begin{enumerate}
\item With the demo project open, click any block of \file{deposit.csv}. The record panel shows its whole row: \col{XC}, \col{YC}, \col{ZC}, \col{FE}, \col{SG}, \col{LITO}.
\item Press \key{i}. File info lists the block size and counts micro inferred.
\item Press \key{⇧B}. Every block is outlined.
\item Zoom in on the cut face until you can count the blocks along a bench.
\end{enumerate}
\end{tryit}

\section{Sub-blocked models}\label{sec:blocks-sub}

A\index{sub-blocked model}\index{block model!sub-blocked} sub-blocked model refines its parent blocks into smaller ones near a contact or a high-grade zone. micro reads these and draws each block at its \emph{true size}, straight from the file.

\begin{figure}[htbp]\centering
\begin{tikzpicture}[x=9mm, y=9mm, font=\sffamily\footnotesize]
\foreach \i in {0,1,2,3} \foreach \j in {0,1} \draw[rule, fill=tint] (\i,\j) rectangle ++(1,1);
\foreach \i in {0,1} \foreach \j in {0,1} \draw[rule, fill=ambersoft] (1+0.5*\i,0.5*\j) rectangle ++(0.5,0.5);
\foreach \i in {0,1} \foreach \j in {0,1} \draw[rule, fill=ambersoft] (2+0.5*\i,1+0.5*\j) rectangle ++(0.5,0.5);
\foreach \i in {0,1} \foreach \j in {0,1} \draw[rule, fill=amberlight!60] (2+0.25*\i,1+0.25*\j) rectangle ++(0.25,0.25);
\draw[amber, line width=0.9pt] (0.6,-0.1) .. controls (1.6,0.4) and (2.0,1.2) .. (3.4,2.1);
\node[text=muted, anchor=west] at (4.3,1.75) {parent 20 m};
\node[text=muted, anchor=west] at (4.3,1.25) {sub-block 10 m};
\node[text=muted, anchor=west] at (4.3,0.75) {sub-block 5 m};
\node[text=amber, anchor=west] at (4.3,0.25) {the contact};
\end{tikzpicture}
\caption{A sub-blocked model in section. micro places every block on one fine lattice (here at half the smallest block size) and gives each its own size in steps of that lattice, so a parent and its sub-blocks are each one record, drawn and weighed at their real volume.}\label{fig:blocks-subblocks}
\end{figure}

\shot{blocks-subblocked}{A sub-blocked model in micro: 20~m parent blocks away from a dipping contact, refined to 10~m and 5~m boxes along it, each drawn at its own size, with block edges on. The status line reads \ui{9,744 sub-blocks · fine 2.5×2.5×2.5 · 3 shapes · largest 20×20×20}.}

Detection is automatic, from the columns that carry each block's size:

\begin{itemize}
\item In a CSV or Parquet model, three per-block dimension columns: \col{DX} \col{DY} \col{DZ}, \col{XINC} \col{YINC} \col{ZINC}, or \col{DIMX} \col{DIMY} \col{DIMZ}.
\item In a Datamine \file{.dm} model, \col{XINC} \col{YINC} \col{ZINC} that vary from record to record.
\end{itemize}

When those are present, micro builds a fine lattice whose pitch on each axis is half the largest step every block size there is a whole number of: half the smallest size when the smaller sizes divide the larger, as they usually do, and half of 1~m for a model with 3~m and 4~m blocks. That is fine enough that the centroid of every block, parent or sub-block, lands on one of its nodes (\cref{fig:blocks-subblocks}). Each block then carries its size as three small whole numbers, its length along each axis in steps of that lattice, so a sub-blocked model costs about a quarter more graphics memory than a regular one, whatever the number of shapes. The status line reports the lattice and the shapes, as in \ui{fine 2.5×2.5×2.5 · 3 shapes · largest 20×20×20} for a model of 20, 10 and 5~m cubes, and file info lists the sizes.

A block's volume is its size in steps times the pitch on each axis: the box drawn on screen. Every analysis that weighs blocks (grade–tonnage, swath, the cross-tab, the resource table, statistics weighted by volume, the reconciliation factors, the contact profile) takes it from there, so they all weigh a block the same way.

Finding the fine lattice means reading every block's centroid and sizes before the first box is drawn, so a sub-blocked \file{.dm} file is read twice on opening. Both reads run on background workers, a run of the file's pages each, and the answers are merged in file order, so the lattice, the sizes and the record numbers are the ones a single reader would find (\cref{sec:opening-resident}).

Everything that works record by record treats a sub-block as an ordinary block: filter, color, pick, statistics, export. The analyses that weigh blocks use each block's own volume, so grade–tonnage and swath do not mistake a 5~m sub-block for a 20~m parent. Join and reconcile aggregate sub-blocks onto a regular grid by the volume each one overlaps (\cref{ch:joins}). Whether a model's sub-blocks actually fill their parents — no gaps, no overlaps, none reaching out of its parent — is a question for Model QA (\cref{sec:qa-subblocks}).

\begin{caution}
Any whole-number refinement works, odd ratios included, and there is no limit on the number of distinct shapes (a real Isatis export has more than three thousand). The one limit is length: a block may span at most 255 steps of the fine lattice along an axis, so a model whose largest block is more than about 127 times its size step opens as points, and the status line says which axis and why. The grades are all there; only the boxes are missing.
\end{caution}

\section{Rotated models}\label{sec:blocks-rotated}

A\index{rotated model}\index{block model!rotated} deposit modeled along its strike has lattice axes that are not east, north and up. Its centroids sit on no axis-aligned grid, so the ordinary search fails, and micro then tries a rotated one. If every centroid sits on a rotated lattice, micro tells you and asks:

\begin{quote}
\sffamily\small\color{ink}deposit.csv isn't on an axis-aligned grid, but every centroid sits on a ROTATED lattice: … Draw it as rotated boxes, or keep the centroids as points?
\end{quote}

The buttons are \ui{rotated boxes} and \ui{points}. Choose \ui{rotated boxes} and the model draws as true rotated boxes, with exact picking and true sections; the status line gives the azimuth and, for a model tilted out of the horizontal, the tilt.

\begin{why}
An inferred rotation is never applied silently. A rotation a hundredth of a degree off moves a block two kilometers along the model by a third of a meter, and a wrong one looks entirely plausible on screen. So micro shows what it found and you decide. The answer is recorded in the layer's lineage, where a reviewer can read it.
\end{why}

\subsection{Declared lattices}

A file can also \emph{state} its lattice instead of leaving it to be inferred. A Parquet file may carry a \code{micro:model} entry in its metadata giving the origin, the three axis directions and the counts; a CSV may carry the same description as a comment line at its top, \code{\# micro:model \{…\}}. micro writes one into every block-model export, an OMF project supplies one, and any producer can (\cref{app:host}). A declared lattice is applied without asking, but it is checked against every centroid. If a block falls off the stated lattice, micro refuses the lattice, says so on the status line, and opens the centroids as points.

\index{rotated model!export}A rotated or sub-blocked model therefore leaves micro \emph{with} its lattice, in either format: exported to CSV or Parquet, it reopens as the same boxes in the same place, declared rather than inferred. That holds for a subset too. An export of the filter's matches or of the selection writes the lattice alone, without the whole model's extent and categories, which would describe rows the file does not hold; micro finds those again on reopening, and the subset opens as boxes of the same lattice, only those blocks.

\subsection{What works on a rotated model}

Everything that works record by record works on a rotated model: filters, statistics, grade–tonnage, swath, the cross-tab, export, and select, flag or proportions by a solid or surface, sub-blocked or not (\cref{ch:geometry}). Join and reconcile work between two models that share the same rotated lattice, a sub-blocked model onto its rotated parent model included.

Two things are not available yet, and micro refuses each with a reason instead of approximating: resampling between models of different orientation, and the empty-model creator. Grade shells, model QA's lattice checks and block vs sample work on a rotated model as on an axis-aligned one.

\subsection{Rotated Datamine models}\label{sec:blocks-dm-rotated}

\index{Datamine .dm!rotated}\index{rotated model!Datamine}
A Datamine model states its rotation in its header: up to three angles \col{ANGLE1}–\col{ANGLE3}, each about an axis \col{ROTAXIS1}–\col{ROTAXIS3} (1 = X, 2 = Y, 3 = Z). Its \col{XC} \col{YC} \col{ZC} are stored in the model's own, unrotated frame. A model with \emph{one} rotation, \col{ANGLE1} about Z, is placed in world coordinates: \col{ANGLE1} turns the model's axes clockwise seen from above, so \code{ANGLE1 = -50} puts the model's X axis at N40°E, and a block's world position is the model's origin plus the rotated \col{XC} \col{YC} \col{ZC}. The convention was checked on a real export, whose block bounds it reproduces to the centimetre. The status line names the rotation and the origin.

The origin is where the model's local (0, 0, 0) sits in the world. micro takes it from the header's \col{X0} \col{Y0} \col{Z0}. A file that carries none, as an anonymised export does, opens about (0, 0, 0), and the status line says so. Type the origin in the open dialog, where a rotated \file{.dm} row shows \ui{origin X Y Z} filled in from the file, or set it later with \menu{Data \& file>Datamine origin…} on the layer's right-click menu (or the command palette). The model moves; the origin is saved with the project.

From then on the coordinate \emph{columns} mean what they mean for any other model: \col{XC} \col{YC} \col{ZC} read in world coordinates, everywhere. A filter \code{XC > 432900}, a swath, statistics, a selection by solid, the record panel, an export: all of them see the world value. The record panel also lists the model's own coordinates beside them, as \ui{XC · model frame} and so on. Volumes are unchanged, since a rotation keeps them.

\index{Datamine .dm!XL YL ZL}The model's own frame is also there as columns: micro gives a placed model three calculated columns, \col{XL} \col{YL} \col{ZL}, each block's centroid in the model's frame — the values the file holds in \col{XC} \col{YC} \col{ZC}. They are ordinary ƒ columns (\cref{ch:columns}), the rotation and the origin written out in their formulas, so a filter \code{XL > 600}, a swath along the model's own axis, statistics or an export can use them like any column. A new origin rewrites the formulas; the values do not move. A file that already has a column of one of those names keeps its own, and an ƒ column of that name you wrote yourself is left alone.

An export writes the world \col{XC} \col{YC} \col{ZC} with the sizes, and the rotated lattice with them: in the \code{micro:model} entry of a Parquet file, in the \code{\# micro:model} line at the top of a CSV. Either reopens as rotated boxes, sub-blocked or not, a filtered subset included (\cref{sec:blocks-rotated}).

\index{Datamine .dm!placement columns}The header's placement constants are \emph{not} exported: \col{ANGLE1}–\col{ANGLE3}, \col{ROTAXIS1}–\col{ROTAXIS3}, \col{X0} \col{Y0} \col{Z0} and \col{XMORIG} \col{YMORIG} \col{ZMORIG}. They describe the model's own frame, and the exported coordinates are world coordinates; a package that read \col{ANGLE1} beside them and applied it would rotate the model twice. \ui{Export rows…} lists them greyed out, \emph{not exported}, with the reason. The counts \col{NX} \col{NY} \col{NZ} and the sizes, which mean the same in any frame, stay.

\begin{tip}
For a package that reads neither \code{micro:model} carrier, tick \ui{lattice as columns} in \ui{Export rows…}. It writes the lattice the Datamine way, as constant columns on every row, consistent with each other: \col{X0} \col{Y0} \col{Z0} (the world position of the model's origin), \col{ANGLE1} with \col{ROTAXIS1} = 3, and \col{XL} \col{YL} \col{ZL}, each block's centroid in the model's own frame, so that world = \col{X0}… + the \col{ANGLE1} rotation of \col{XL}…; a regular model adds \col{XMORIG}…, \col{XINC}… and \col{NX}…. It costs every row (the dialog says how much); micro itself reads the lattice from the comment line, not from these columns. A lattice tilted out of the horizontal gets no \col{ANGLE1} (Datamine's three-rotation order is unverified), and a sub-blocked model no parent lattice, which micro does not know.
\end{tip}

A model with two or three rotations is still drawn in its own, unrotated frame, with a warning on the status line: the order in which Datamine composes them has not been verified on a real file, and a placement that looks plausible and is wrong is worse than an honest frame. Its volumes, grades and tonnages are right; its positions against other layers are not.

\section{Datamine \file{.dm} models}\label{sec:blocks-dm}

A\index{Datamine .dm} \file{.dm} file stores fixed-width records, so a given field always sits at the same offset inside every record. micro uses that to read one column without touching the others: coloring by \col{FE} reads \col{FE}, a filter on \col{LITO} reads \col{LITO}. On a model with many columns this is several times faster than decoding every record.

The header also names the columns and types them, so there is nothing to sniff and no mapping to confirm. A regular model states its lattice in the header as constants (\col{XMORIG}…, \col{XINC}…, \col{NX}…); sub-blocking is detected from the per-record \col{XINC} \col{YINC} \col{ZINC} (\cref{sec:blocks-sub}). A file with neither, whose per-record sizes are all the same or which has none, has its lattice found from the centroids, as a CSV has, checked against the stated sizes; a rotated lattice found that way is offered, never applied silently (\cref{sec:blocks-rotated}).

micro reads \file{.dm} and does not write it. To hand a model on, export it as CSV or Parquet (\cref{sec:columns-export}), which also makes micro a \file{.dm} to CSV converter. A \file{.dm} file can instead be half of a \emph{wireframe} (\cref{sec:meshes-dm}).

\section{Parquet models}\label{sec:blocks-parquet}

Parquet\index{Parquet!block models} is a columnar file format: values are stored column by column, compressed, in blocks of rows called \term{row groups}, with a footer that lists every column's minimum, maximum and blank count for each group. Three properties make it a good home for a block model.

\textbf{It opens from its footer.} The open dialog's description (rows, columns, compression) is read from the last few kilobytes of the file. The model then streams one row group at a time, and nothing decoded is held.

\textbf{Filters skip what cannot match.} For a filter such as \code{FE > 60}, micro reads each row group's maximum of \col{FE} from the footer and skips every group whose maximum is below 60, without reading it. It also reads only the columns the expression names. On a model stored by micro the same trick works in space: a section, a selection box or a solid skips the row groups whose bounding box lies outside it.

\textbf{It is typed and self-describing.} A number is a number; blanks are blanks, not a code. A model micro wrote also records its grid, its extents and its categories in the file, so it reopens without the inference pass of \cref{sec:blocks-lattice}.

\subsection{Store as Parquet…}\label{sec:blocks-optimize}

\index{Parquet!Store as Parquet}\index{optimize to spatial Parquet}
Right-click a CSV or \file{.dm} model and choose \menu{Data \& file>Store as Parquet…}, or take the offer micro makes after a model has been slow to open (\cref{sec:opening-resident}). The dialog, \ui{store block models as Parquet}, lists the model with its size and block count, and what will happen to it: \ui{→ Morton-sorted (spatial index)}.

\shot[0.75\linewidth]{blocks-store-parquet}{The \ui{store block models as Parquet} dialog for the demo deposit: what will be converted, how, and whether the original stays.}

\term{Morton order} \parencite{morton1966} numbers blocks so that blocks close in space are close in the file. Sorted that way, each row group covers a compact piece of the deposit, and its footer entry, the minimum and maximum of X, Y and Z, becomes an entry in a spatial index. That is what makes the spatial skipping above work.

\begin{table}[htbp]
\small
\begin{tabularx}{\linewidth}{@{}NND@{}}\toprule
setting & meaning & default \\ \midrule
the checkbox on each row & convert this model & on \\
\ui{keep the original file in the project} & leave the CSV or \file{.dm} where it is; otherwise it is removed after the Parquet has saved & off \\
\ui{do this automatically on every save} & shown when the dialog was raised by a project save; makes the conversion standing, for models under 8 million blocks & off \\
\bottomrule
\end{tabularx}
\caption{The \ui{Store as Parquet…} dialog. \ui{Convert} runs it; \ui{Not now} leaves the model as it is.}\label{tab:blocks-optimize}
\end{table}

The layer swaps to the Parquet file in place. Its styling and filter are re-applied. Painted columns are reordered along with the blocks and stay editable. The layer's lineage records the conversion.

The same conversion is offered when you save a project. Tick \ui{do this automatically on every save} to make it standing; \ui{Not now} asks again in your next session, and a conversion without the tick still asks about the next new model. \menu{Tools>Auto-optimize block models to spatial Parquet on save} shows the standing choice and changes it, and \menu{Tools>Settings…} can set it to never. Models above 8 million blocks are left for the explicit menu command: they sort in bounded memory, so they do not fail, but rewriting a multi-gigabyte file in the middle of a save would be an unwelcome surprise.

\begin{note}
\ui{Export rows…} can also write Parquet, with a \ui{spatial sort} checkbox for the same Morton order (\cref{sec:columns-export}). Use \ui{Store as Parquet…} to change what the \emph{project} keeps, and \ui{Export rows…} to make a file for someone else.
\end{note}

\subsection{Build index}

A model you would rather keep as CSV or \file{.dm} can still be made quick to filter. \menu{Data \& file>Build index (stats + pushdown)…} on its menu streams the model once and records, for consecutive stretches of the file, each column's range. A later filter then skips the stretches that cannot match, as Parquet's footer allows. The index is written as a small file beside the model on the next project save (\cref{ch:store}); delete it and micro scans again.

\section{Which format to keep a model in}\label{sec:blocks-which}

Each of the three formats wins somewhere (\cref{tab:blocks-formats}).

\begin{table}[htbp]
\small
\begin{tabularx}{\linewidth}{@{}lLL@{}}\toprule
format & good at & weak at \\ \midrule
CSV & read by everything; readable by eye; the lingua franca between packages & largest on disk; every number is parsed from text; types and blanks are guessed \\
Datamine \file{.dm} & opens at once; one column read at a time; typed & micro cannot write it; a model with two or three rotations is not placed in world coordinates \\
Parquet & smallest; typed; filters skip row groups, and regions of space once Morton-sorted; reopens without inference & not readable by eye; not every mining package reads it \\
\bottomrule
\end{tabularx}
\caption{The three block-model formats compared.}\label{tab:blocks-formats}
\end{table}

A workable rule: \emph{receive} a model in whatever format its author's software writes, \emph{work} on it as Parquet, and \emph{deliver} CSV when the recipient needs it. Accepting the save-time offer does the middle step for you.

Keep the original when it is the record of what you were sent. A project is a plain folder (\cref{ch:projects}); a model and the Parquet micro made from it can sit side by side, and the lineage of the Parquet layer names its source.

\section{Grade shells}\label{sec:blocks-shells}

A \term{grade shell} is the surface that encloses every block at or above a cutoff: the three-dimensional equivalent of a contour line. It answers ``where is the material above 55\,\% Fe, as a shape?'', and it gives you that shape as a solid you can measure and reuse.

Right-click a block model and choose \menu{Analyse>Grade shell…} (or press \key{⌘K} and type \code{shell}). The window has four settings and one button.

\begin{table}[htbp]
\small
\begin{tabularx}{\linewidth}{@{}lMD@{}}\toprule
setting & meaning & default \\ \midrule
\ui{column} & the numeric column to contour & the column the layer is colored by \\
\ui{cutoff (≥)} & the shell encloses every block at or above this value & none — type one \\
\ui{smoothing} & average each block with its neighbors before contouring: \ui{none (the blocks as they are)}, \ui{light}, \ui{medium}, \ui{strong} & none \\
\ui{only the filtered blocks} & count only blocks matching the layer's filter as inside; available when the layer has a filter & off \\
\bottomrule
\end{tabularx}
\caption{The grade shell window. \ui{Make shell} builds the mesh.}\label{tab:blocks-shell}
\end{table}

\ui{Make shell} reads the model and adds a new mesh layer named after what it is, such as \ui{FE ≥ 55 shell}, or \ui{FE ≥ 55 shell (smoothed ×1)} when smoothing was used. The window reports the triangle count and how many blocks ended up inside, as in \ui{62,232 triangles around 11,029 blocks}. Each click makes a new layer, so shells of different settings can be compared side by side.

The surface passes between block centers where the grade crosses the cutoff; micro traces it by marching tetrahedra \parencite{doi1991}. Where the model ends, or a block is missing, it closes on the block face. So a shell is always a closed solid, never an open sheet. With smoothing off it shows exactly the blockiness the model has, which is honest and a little ugly; \ui{light} is usually enough to make a shape that reads as geology without moving it far.

\ui{only the filtered blocks} is how you make a shell within one domain: filter the model to \code{LITO = "HEMATITE"}, tick the box, and blocks of other lithologies count as outside whatever their grade.

\begin{tryit}
Make a high-grade shell in the demo.
\begin{enumerate}
\item Right-click \file{deposit.csv} and choose \menu{Analyse>Grade shell…}.
\item Leave \ui{column} on \col{FE}, type \code{60} for the cutoff, and click \ui{Make shell}.
\item The window reports \ui{62,232 triangles around 11,029 blocks}: every block at or above 60\,\% Fe, which is the 11,013 of the demo's \code{HEMATITE} core and 16 itabirite blocks. A layer \ui{FE ≥ 60 shell} appears. Hide \file{deposit.csv} with its eye (or press \key{h} with it active) to see the shell by itself.
\item Set \ui{smoothing} to \ui{light} and click \ui{Make shell} again. A second layer, \ui{FE ≥ 60 shell (smoothed ×1)}, appears with about half the triangles (32,212). Compare the two shells by toggling their eyes.
\item Right-click a shell and choose \menu{Analyse>Mesh QA / volume…}: it is watertight, and it has a volume.
\end{enumerate}
\end{tryit}

\shot{blocks-shell}{A grade shell on the demo deposit: every block at or above 60\,\% Fe, lightly smoothed and enclosed in a closed mesh, with the model hidden and the four fans of drillholes threading it.}

A shell is a mesh like any wireframe. You can select or flag the model, or the drillholes, against it (\cref{ch:geometry}), read its volume and tonnage in mesh QA (\cref{ch:qa}), restyle it, and export it (\cref{sec:meshes-export}). Its lineage records the column, the cutoff, the smoothing and, when used, the filter.

Shells are built for regular models, axis-aligned or rotated: a rotated model is read and contoured in its own lattice frame and the shell is turned into world. On a sub-blocked model the command says why it cannot run.

\subsection{In a routine}

The tool is \code{shell}. A re-run replaces the shell an earlier run made instead of adding a copy.

\begin{lstlisting}[language=yaml]
- tool: "shell"
  with:
    layer: "deposit.csv"
    col: "FE"
    cutoff: 60
    smooth: 1
\end{lstlisting}

\code{smooth} is the number of smoothing passes, 0 to 3, matching \ui{none} to \ui{strong}. An optional \code{name} sets the mesh's name; without it the shell is named as the window names it. Routines are the subject of \cref{ch:routines}.

\section{Related tools}\label{sec:blocks-qa}\label{sec:blocks-new}

Two tools that belong to block models are described elsewhere.

\textbf{Model QA.} Before anyone reports from a model, run \menu{Analyse>Model QA…} from its menu. It reads the model once and lists what a reviewer should see: duplicate blocks, centroids off the grid, gaps inside a column of blocks, blanks, sentinel codes such as \code{-99} read as numbers, densities outside a plausible range, and outliers. Click a finding to select its blocks; the report changes nothing in the model (\cref{ch:qa}).

\textbf{Empty models and common grids.} \menu{Data>New grid / block model…} makes a model that does not exist yet: a blank lattice to paint domains into, or one that two models of different block sizes can both be compared on. It works from typed parameters, from the extent of a layer, or by covering several compatible models at once (\cref{sec:grids-new}; the reconciliation it serves is in \cref{ch:joins}).
